\label{sec:itemfeat-text}
\begin{table}[H]
\centering
\caption{Item-difficulty transfer: within-benchmark Spearman correlation and pairwise accuracy between predicted easiness and observed pass rate, averaged across benchmarks. ID: items held out within the 11 training benchmarks; OOD validation and OOD test: three unseen benchmarks each. The additional pass-rate and Rasch-target probes use benchmark-grouped selection and were not evaluated on ID (--).}
\label{tab:item-difficulty}
\small\setlength{\tabcolsep}{4pt}
\begin{tabular}{lcccccc}
\toprule
 & \multicolumn{2}{c}{ID} & \multicolumn{2}{c}{OOD validation} & \multicolumn{2}{c}{OOD test}\\
\cmidrule(lr){2-3}\cmidrule(lr){4-5}\cmidrule(lr){6-7}
Probe & Spearman & Pairwise & Spearman & Pairwise & Spearman & Pairwise\\
\midrule
Random score & $-0.003$ & $0.499$ & $-0.024$ & $0.490$ & $0.076$ & $0.529$ \\
Log character count & $0.107$ & $0.539$ & $0.058$ & $0.525$ & $0.061$ & $0.521$ \\
Ridge on embeddings & $0.400$ & $0.649$ & $0.196$ & $0.570$ & $-0.020$ & $0.494$ \\
Pass-rate ridge & -- & -- & $0.165$ & $0.557$ & $-0.035$ & $0.487$ \\
Rasch-target ridge & -- & -- & $0.130$ & $0.545$ & $-0.035$ & $0.487$ \\
Pass-rate MLP & -- & -- & $0.071$ & $0.523$ & $0.109$ & $0.539$ \\
Rasch-target MLP & -- & -- & $0.120$ & $0.540$ & $0.104$ & $0.538$ \\
\bottomrule
\end{tabular}

\end{table}

\begin{qnote}{Q2}Section 5.4: the probe table and the prior-mean table now sit in Results, with the probes defined in Section 4.2. Does the abstract say no more than these two tables show?\end{qnote}
Embeddings show a clear signal on known benchmarks. On the six unseen benchmarks the two MLP probes pass the criterion we fixed in advance, with benchmark-macro Spearman $0.090$ and $0.103$ against $0.026$ for a random baseline and a positive correlation on 5 of 6 benchmarks; the ridge probes do not ($0.065$ and $0.047$). On the three test benchmarks alone, the pass-rate MLP reaches $0.109$, while the Rasch-target MLP reaches $0.059$, below the sampled random control's $0.076$. The aggregate rank criterion therefore does not establish a consistent test-only advantage.

The rank signal yields only a small gain at the validation-selected scale. In Table~\ref{tab:probe-prior}, every $\gamma\ge0.25$ is worse than $\gamma=0$ on both splits. The scale selected on validation ($\gamma=0.05$, Rasch-target MLP) improves Brier ALC by about $0.0004$ on each split in its matched replay; this comparison alone does not establish a reliable improvement. This tests these two probes only; a stronger difficulty predictor, for instance one trained through an item response likelihood on more items, could do better.

\begin{table}[H]
\centering
\caption{Predicted item difficulty as the prior mean of the Rasch item deviation, Equation~\eqref{eq:probe-prior}. Cells are $\Delta_{m-\mathrm{J}}$ of the fitted Rasch with scale $\gamma$, competition-shaped draws; $\gamma=0$ is the fitted Rasch without a prior mean. Bold: lowest in each column.}
\label{tab:probe-prior}
\begin{tabular}{lcccc}
\toprule
 & \multicolumn{2}{c}{Pass-rate MLP} & \multicolumn{2}{c}{Rasch-target MLP} \\
Scale $\gamma$ & OOD validation & OOD test & OOD validation & OOD test \\
\midrule
$0$ & $\mathbf{+0.0052}$ & $\minus0.0119$ & $+0.0052$ & $\minus0.0119$ \\
\midrule
$0.05$ & $+0.0062$ & $\mathbf{\minus0.0119}$ & $\mathbf{+0.0048}$ & $\mathbf{\minus0.0123}$ \\
$0.1$ & $+0.0103$ & $\minus0.0082$ & $+0.0051$ & $\minus0.0120$ \\
$0.25$ & $+0.0332$ & $+0.0150$ & $+0.0099$ & $\minus0.0079$ \\
$0.5$ & $+0.0690$ & $+0.0491$ & $+0.0273$ & $+0.0069$ \\
$1$ & $+0.1065$ & $+0.0826$ & $+0.0704$ & $+0.0433$ \\
$1.5$ & $+0.1326$ & $+0.1050$ & $+0.1040$ & $+0.0724$ \\
\bottomrule
\end{tabular}

\end{table}

